\chapter{Implementatie}

Dit hoofdstuk beschrijft de concrete uitwerking van het proof of concept dat werd ontwikkeld in het kader van dit onderzoek. Zoals uiteengezet in de methodologie, werd de implementatie stapsgewijs opgebouwd: eerst een werkende basisstructuur, vervolgens de integratie van de AR-laag. Dit hoofdstuk volgt die logische opbouw en documenteert per component de technische keuzes, de implementatiedetails en de aansluiting op de in de literatuurstudie beschreven vereisten.

\section{Technische omgeving en tools}

De applicatie werd ontwikkeld als een native iOS-applicatie in de programmeertaal Swift (versie 5.9), binnen de ontwikkelomgeving Xcode 15. De keuze voor een native iOS-oplossing is gemotiveerd door de rijke ondersteuning van Apple voor AR-toepassingen via het ARKit-framework, de stabiele integratie met de hardware van de iPhone (camera, GPS, gyroscoop, versnellingsmeter) en de betere prestatiegaranties in vergelijking met cross-platform alternatieven.

De voornaamste frameworks die gebruikt worden, zijn:

\begin{itemize}
    \item \textbf{ARKit} – Apple’s framework voor augmented reality op iOS.
    \item \textbf{SceneKit} – 3D-renderingframework voor AR-objecten.
    \item \textbf{MapKit} – kaartweergave en annotaties.
    \item \textbf{CoreLocation} – GPS en kompasrichting.
    \item \textbf{Combine} – reactief programmeerframework.
\end{itemize}

De applicatie maakt gebruik van het SwiftUI app lifecycle-model, gecombineerd met UIKit via een bridge.

\section{Softwarearchitectuur}

De architectuur volgt het principe van \textit{separation of concerns} en het Repository-patroon.

\subsection{Datamodellen}

Het centrale model is \texttt{GraveRecord}, een structuur met grafgegevens zoals naam, coördinaten en metadata.

Daarnaast wordt \texttt{ARNavigationState} gebruikt als \texttt{ObservableObject} voor de navigatiestatus.

De bearing-berekening gebeurt via:

\begin{verbatim}
    y = sin(Δlon) × cos(lat2)
    x = cos(lat1) × sin(lat2) − sin(lat1) × cos(lat2) × cos(Δlon)
    bearing = atan2(y, x)
\end{verbatim}

\subsection{Dataservice}

De dataservice volgt het Repository-patroon en leest JSON-data in via asynchrone verwerking. Zoekfunctionaliteit gebeurt via case-insensitieve matching.

\subsection{Locatieservice}

De LocationManager levert GPS- en headingdata met hoge nauwkeurigheid en werkt met delegates.

\section{De AR-component}

\subsection{ARKit-configuratie}

De applicatie gebruikt \texttt{ARWorldTrackingConfiguration} met:

\begin{itemize}
    \item worldAlignment = gravityAndHeading
    \item horizontale vlakdetectie
\end{itemize}

\subsection{GPS naar AR-coördinaten}

De omzetting gebeurt via:

\[
X = \Delta lon \times 111.320 \times \cos(lat) \times 1000
\]
\[
Z = -\Delta lat \times 111.320 \times 1000
\]

\subsection{Navigatiepijl}

De pijl wordt als 3D-object opgebouwd en per frame geüpdatet.

Rotatie gebeurt via:

\[
\theta = \text{atan2}(dx, dz)
\]

\subsection{Grafmarker}

Bestaat uit:

\begin{itemize}
    \item Roterende kubus
    \item Verticale lichtbalk
    \item Pulserende grondring
    \item Billboard-label
\end{itemize}

\section{UI-componenten}

\subsection{ViewController}

Centrale controller met volgende lagen:

\begin{enumerate}
    \item AR-weergave
    \item Minikaart
    \item Info-paneel
    \item Zoekknop
    \item Statuslabels
\end{enumerate}

\subsection{Minikaart}

MapKit-view met fullscreen animatie en annotaties.

\subsection{Zoekscherm}

Geïmplementeerd via \texttt{UISheetPresentationController} met live zoekresultaten.

\subsection{Grafinformatie}

Blur-card met grafdetails en sluitfunctie.

\section{Mockdata}

Data wordt opgeslagen in JSON-formaat met velden zoals naam, coördinaten en sectie.

\section{Iteratief ontwikkelproces}

\begin{itemize}
    \item Iteratie 1: basisstructuur
    \item Iteratie 2: zoekfunctionaliteit
    \item Iteratie 3: AR-ankers
    \item Iteratie 4: 2D-pijl
    \item Iteratie 5: 3D-pijl
    \item Iteratie 6: UI-verfijning
\end{itemize}

\section{Technische beperkingen}

\begin{itemize}
    \item GPS-nauwkeurigheid (3–10 meter)
    \item Initialisatietijd
    \item AR drift
    \item Licht- en weersinvloeden
\end{itemize}

\section{Aansluiting op methodologie en literatuur}

De implementatie ondersteunt zowel kaart- als AR-navigatie en sluit aan bij de onderzoeksvraag.

Belangrijke literatuurinzichten:

\begin{itemize}
    \item Verminderde cognitieve belasting door AR
    \item Directe visuele navigatie
    \item Gebruiksvriendelijke interface
\end{itemize}